\documentclass[11pt,letterpaper]{article}
% KJO_preamble.tex
% Shared LaTeX preamble for all KJO documentation documents.
% \input this file immediately after \documentclass{}.

% --- Encoding and fonts -------------------------------------------------------
\usepackage[utf8]{inputenc}
\usepackage[T1]{fontenc}
\usepackage{lmodern}
\usepackage{microtype}

% --- Page geometry ------------------------------------------------------------
\usepackage[
  letterpaper,
  top=1in, bottom=1in,
  left=1.25in, right=1in
]{geometry}

% --- Core utilities -----------------------------------------------------------
\usepackage{amsmath}
\usepackage{graphicx}
\usepackage{xcolor}
\usepackage{parskip}        % space between paragraphs; no indent
\usepackage{enumitem}
\usepackage{caption}
\usepackage{subcaption}
\usepackage{float}

% --- Tables -------------------------------------------------------------------
\usepackage{booktabs}
\usepackage{longtable}
\usepackage{array}
\usepackage{multirow}

% --- Code listings ------------------------------------------------------------
\usepackage{listings}
\lstset{
  basicstyle=\ttfamily\small,
  keywordstyle=\bfseries\color{KJOblue},
  commentstyle=\itshape\color{gray},
  stringstyle=\color{KJOgreen!70!black},
  breaklines=true,
  frame=single,
  framesep=4pt,
  xleftmargin=6pt,
  xrightmargin=4pt,
  numbers=left,
  numberstyle=\tiny\color{gray},
  numbersep=8pt,
  language=C++,
}

% --- TikZ / PGF ---------------------------------------------------------------
\usepackage{tikz}
\usetikzlibrary{
  arrows.meta,
  shapes.geometric,
  shapes.misc,
  shapes.symbols,
  positioning,
  fit,
  calc,
  backgrounds,
  decorations.pathreplacing,
  decorations.markings,
  matrix,
}
\usepackage{pgfplots}
\pgfplotsset{compat=1.18}

% --- Headers / footers --------------------------------------------------------
\usepackage{fancyhdr}
\usepackage{lastpage}
\pagestyle{fancy}
\fancyhf{}
\renewcommand{\headrulewidth}{0.4pt}
\renewcommand{\footrulewidth}{0.4pt}
% Per-document: \lhead{}, \rhead{}, \cfoot{} are set in the main .tex file.

% --- Section formatting -------------------------------------------------------
\usepackage{titlesec}
\titleformat{\section}{\large\bfseries\color{KJOblue}}{{\thesection}}{0.8em}{}[\titlerule]
\titleformat{\subsection}{\normalsize\bfseries\color{KJOblue!80!black}}{{\thesubsection}}{0.8em}{}
\titleformat{\subsubsection}{\normalsize\bfseries}{{\thesubsubsection}}{0.8em}{}

% --- Hyperlinks ---------------------------------------------------------------
\usepackage[
  colorlinks=true,
  linkcolor=KJOblue,
  urlcolor=KJOblue!70!black,
  citecolor=KJOblue,
  pdfborder={0 0 0},
]{hyperref}

% --- Color palette -----------------------------------------------------------
\definecolor{KJOblue}   {RGB}{ 30,  90, 160}   % RCU
\definecolor{KJOred}    {RGB}{180,  30,  30}   % EMU
\definecolor{KJOgreen}  {RGB}{ 20, 120,  60}   % GSEMU
\definecolor{KJOorange} {RGB}{200, 110,  20}   % GSE / pressurization
\definecolor{KJOgray}   {RGB}{100, 100, 100}   % Buses / links
\definecolor{KJOpurple} {RGB}{110,  40, 140}   % Shared libraries
\definecolor{KJObg}     {RGB}{248, 248, 252}   % Light background for code

% --- Convenience macros -------------------------------------------------------

% Hardware component name (small caps)
\newcommand{\hw}[1]{\textsc{#1}}

% Software module name (monospace)
\newcommand{\sw}[1]{\texttt{#1}}

% Command name (monospace bold)
\newcommand{\cmd}[1]{\texttt{\textbf{#1}}}

% Pin / signal name
\newcommand{\sig}[1]{\texttt{#1}}

% Note / callout box
\usepackage{tcolorbox}
\tcbuselibrary{skins}
\newtcolorbox{notebox}{
  enhanced,
  colback=KJOblue!5,
  colframe=KJOblue!40,
  fonttitle=\bfseries,
  title=Note,
  left=4pt, right=4pt, top=2pt, bottom=2pt,
}
\newtcolorbox{warningbox}{
  enhanced,
  colback=KJOred!5,
  colframe=KJOred!40,
  fonttitle=\bfseries\color{KJOred},
  title=Warning,
  left=4pt, right=4pt, top=2pt, bottom=2pt,
}

% Title page macro: \KJOtitlepage{title}{subtitle}{version}{date}{description}
\newcommand{\KJOtitlepage}[5]{%
  \begin{titlepage}
    \centering
    \vspace*{2cm}
    {\Huge\bfseries\color{KJOblue} #1 \par}
    \vspace{0.5cm}
    {\Large\color{KJOgray} #2 \par}
    \vspace{1.5cm}
    \rule{0.5\textwidth}{0.4pt}\par
    \vspace{1cm}
    {\large Ken Overton \par}
    \vspace{0.3cm}
    {\normalsize \textit{KJO Liquid Rocket Motor Control System} \par}
    \vspace{1.5cm}
    \begin{tabular}{ll}
      \textbf{Version:} & #3 \\[4pt]
      \textbf{Date:}    & #4 \\
    \end{tabular}
    \vspace{1.5cm}\par
    \begin{minipage}{0.75\textwidth}
      \centering
      \normalsize #5
    \end{minipage}
    \vfill
    {\small\color{KJOgray} \textcopyright\ 2025--2026 Ken Overton.
     All rights reserved. \par}
  \end{titlepage}%
}

\newcommand{\docversion}{0.1}
\newcommand{\docdate}{September 2026}
\lhead{\textsc{KJO Servo Tuner}}
\rhead{\textsc{User Guide} --- v\docversion}
\cfoot{\thepage\ of \pageref{LastPage}}
\hypersetup{pdftitle={KJO Servo Tuner User Guide}, pdfauthor={Ken Overton}}

\begin{document}
% KJO_tikz_styles.tex
% Shared TikZ node and edge styles for all KJO block diagrams.
% \input this file inside the document body before any tikzpicture environments.

\tikzset{
  % ── Hardware block (peripheral chip / PCB component) ──────────────────────
  hwblock/.style={
    rectangle, draw=KJOgray, line width=1pt,
    fill=KJOgray!10, rounded corners=4pt,
    minimum width=2.6cm, minimum height=0.8cm,
    font=\small, align=center, inner sep=4pt,
  },
  % Sub-variant: MCU (larger, blue border)
  mcublock/.style={
    hwblock,
    draw=KJOblue, fill=KJOblue!8,
    minimum width=3.0cm, minimum height=1.0cm,
    font=\small\bfseries,
  },
  % Sub-variant: power / battery
  powerblock/.style={
    hwblock,
    draw=KJOorange, fill=KJOorange!8,
  },
  % ── Software module box ────────────────────────────────────────────────────
  swmodule/.style={
    rectangle, draw=KJOblue, line width=0.8pt,
    fill=KJOblue!6, rounded corners=3pt,
    minimum width=2.8cm, minimum height=0.7cm,
    font=\small, align=center, inner sep=3pt,
  },
  % Sub-variant: shared library
  sharedlib/.style={
    swmodule,
    draw=KJOpurple, fill=KJOpurple!6,
    dashed,
  },
  % ── Bus / link annotation ──────────────────────────────────────────────────
  busline/.style={
    draw=KJOgray, line width=1.5pt,
    -{Stealth[length=6pt,width=4pt]},
  },
  bidbusline/.style={
    draw=KJOgray, line width=1.5pt,
    {Stealth[length=6pt,width=4pt]}-{Stealth[length=6pt,width=4pt]},
  },
  radioline/.style={
    draw=KJOblue!60, line width=1.5pt, dashed,
    {Stealth[length=6pt,width=4pt]}-{Stealth[length=6pt,width=4pt]},
  },
  gpioline/.style={
    draw=KJOgreen!60!black, line width=1.2pt,
    -{Stealth[length=5pt,width=4pt]},
  },
  % ── Flowchart shapes ──────────────────────────────────────────────────────
  flowproc/.style={
    rectangle, draw=KJOblue!70, line width=0.8pt,
    fill=KJOblue!8, rounded corners=3pt,
    minimum width=3.2cm, minimum height=0.7cm,
    font=\small, align=center, inner sep=3pt,
  },
  flowdec/.style={
    diamond, draw=KJOorange, line width=0.8pt,
    fill=KJOorange!8,
    minimum width=2.2cm, minimum height=0.9cm,
    font=\small, align=center, inner sep=1pt,
  },
  flowterminal/.style={
    rounded rectangle, draw=KJOblue, line width=1pt,
    fill=KJOblue!12,
    minimum width=2.8cm, minimum height=0.7cm,
    font=\small\bfseries, align=center, inner sep=3pt,
  },
  flowarrow/.style={
    draw=KJOgray, line width=0.9pt,
    -{Stealth[length=5pt,width=4pt]},
  },
  % ── Propellant schematic styles ────────────────────────────────────────────
  % Generic pipe line
  pipe/.style={
    draw=black, line width=1.8pt,
  },
  pipearrow/.style={
    pipe,
    -{Stealth[length=7pt,width=5pt]},
  },
  % Oxidizer (N2O) — blue
  oxpipe/.style={
    draw=KJOblue, line width=1.8pt,
  },
  oxpipearrow/.style={
    oxpipe,
    -{Stealth[length=7pt,width=5pt]},
  },
  % Fuel — red
  fuelpipe/.style={
    draw=KJOred, line width=1.8pt,
  },
  fuelpipearrow/.style={
    fuelpipe,
    -{Stealth[length=7pt,width=5pt]},
  },
  % GSE fill line — green
  gsepipe/.style={
    draw=KJOgreen!70!black, line width=1.8pt,
  },
  gsepipearrow/.style={
    gsepipe,
    -{Stealth[length=7pt,width=5pt]},
  },
  % Pressurization — orange dashed
  presspipe/.style={
    draw=KJOorange, line width=1.4pt, dashed,
  },
  presspipearrow/.style={
    presspipe,
    -{Stealth[length=6pt,width=4pt]},
  },
  % Tank (tall rectangle, rounded corners)
  tank/.style={
    rectangle, draw=black, line width=1.5pt,
    rounded corners=4pt, fill=white,
    minimum width=1.8cm, minimum height=2.4cm,
    font=\small\bfseries, align=center, inner sep=4pt,
  },
  n2otank/.style={
    tank, draw=KJOblue, fill=KJOblue!8,
  },
  fueltank/.style={
    tank, draw=KJOred, fill=KJOred!8,
  },
  srctank/.style={
    tank, draw=KJOgreen!70!black, fill=KJOgreen!6,
  },
  % Valve box (servo-actuated)
  valvebox/.style={
    rectangle, draw=black, line width=1pt,
    fill=white, minimum width=2.2cm, minimum height=0.65cm,
    font=\footnotesize\bfseries, align=center, inner sep=3pt,
  },
  gseval/.style={
    valvebox, draw=KJOgreen!70!black, fill=KJOgreen!8,
  },
  emuval/.style={
    valvebox, draw=KJOred, fill=KJOred!6,
  },
  % Component (generic inline component)
  component/.style={
    rectangle, draw=KJOgray, line width=0.8pt,
    fill=white!96!gray, minimum width=2.2cm, minimum height=0.6cm,
    font=\footnotesize, align=center, inner sep=3pt,
  },
  % Instrument callout circle (PT, TT, TC)
  instrument/.style={
    circle, draw=KJOgray, line width=0.8pt,
    fill=white, minimum size=0.6cm,
    font=\tiny\bfseries, align=center, inner sep=0pt,
  },
  % Sensor lead line
  sensorlead/.style={
    draw=KJOgray, line width=0.7pt,
  },
  % Boundary box (dashed outline for GSE / Rocket regions)
  gseboundary/.style={
    draw=KJOgreen!50, line width=1pt, dashed,
    rounded corners=6pt,
  },
  rocketboundary/.style={
    draw=KJOred!40, line width=1pt,
    rounded corners=6pt,
  },
  % Sequence diagram lifeline
  lifeline/.style={
    draw=KJOgray, line width=0.6pt, dashed,
  },
  % Sequence diagram message arrow
  seqarrow/.style={
    draw=KJOblue, line width=0.9pt,
    -{Stealth[length=5pt,width=4pt]},
  },
  % State node (FSM)
  statenode/.style={
    circle, draw=KJOblue, line width=1pt,
    fill=KJOblue!10,
    minimum size=1.2cm,
    font=\small\bfseries, align=center,
  },
  statetransition/.style={
    draw=KJOgray, line width=0.9pt,
    -{Stealth[length=5pt,width=4pt]},
    font=\footnotesize,
  },
}


\KJOtitlepage%
  {Servo Tuner}%
  {Bench Tool User Guide}%
  {\docversion}{\docdate}%
  {A standalone bench tool for finding the PWM open/close endpoints of a
   servo-actuated ball valve, using a Serial console to move the servo and
   read back its position potentiometer.}

\tableofcontents\newpage

%═══════════════════════════════════════════════════════════════════════════════
\section{Purpose}
%═══════════════════════════════════════════════════════════════════════════════

\begin{sloppypar}
Servo Tuner is a small standalone sketch used on the bench to determine the
correct \texttt{PWM\_OPEN}/\texttt{PWM\_CLOSE} calibration constants for a
single servo-actuated ball valve, before those values are hardcoded into
production firmware (EMU, GSEMU). It drives one servo channel and reads one
position potentiometer, letting the operator move to a candidate PWM count,
read the resulting pot value, and iteratively nudge the target until the
valve sits exactly at its mechanical open or closed stop.
\end{sloppypar}

This is a different program from the \textbf{EMU Servo Tuner} (a related but
separate tool that tunes all three of EMU's installed valves against its
actual per-valve pin/channel table -- see that tool's own guide).

%═══════════════════════════════════════════════════════════════════════════════
\section{Hardware Required}
%═══════════════════════════════════════════════════════════════════════════════

\begin{itemize}[nosep]
  \item Adafruit Feather M0 (ATSAMD21G18A)
  \item Adafruit 8-channel PWM/Servo FeatherWing (PCA9685), I\textsuperscript{2}C
  \item One servo-actuated ball valve, servo wired to \textbf{PCA9685 channel 7}
  \item The valve's position potentiometer wired to \textbf{analog pin A1}
  \item USB cable for Serial (115200 baud)
\end{itemize}

\begin{notebox}
The servo channel (7) and pot pin (A1) are hardcoded in the sketch, not
configurable at runtime. To tune a valve wired to different pins, edit
\texttt{servo\_channel}/\texttt{pot\_channel} in \texttt{main.cpp} and
reflash.
\end{notebox}

%═══════════════════════════════════════════════════════════════════════════════
\section{How to Use}
%═══════════════════════════════════════════════════════════════════════════════

Open a Serial monitor at 115200 baud after powering the board. The tool
runs an outer ``new target'' loop and an inner ``adjust'' loop:

\begin{enumerate}[nosep]
  \item \textbf{Target}: the tool prints a suggested PWM range (see
        §\ref{sec:suggested-range}) and prompts for a target PWM count.
        Enter a number and press Enter.
  \item \textbf{Hysteresis delta}: enter a count value. Every subsequent
        move first drives the servo \texttt{target~+~delta} counts past
        the target, then back to the target itself -- this removes
        mechanical backlash from the reading, so the pot value reflects
        the valve consistently approaching from one direction.
  \item The servo moves to the target and the tool prints the resulting
        potentiometer reading (\texttt{Encoder = ...}).
  \item \textbf{Action prompt} -- type one letter and press Enter:
        \begin{itemize}[nosep]
          \item \texttt{a} (adjust): enter a signed delta (e.g.\ \texttt{+20}
                or \texttt{-15}); the tool adds it to the current target,
                re-approaches through the hysteresis move, and prints the
                new pot reading.
          \item any other key: returns to step 1 for a completely fresh
                target and hysteresis delta.
        \end{itemize}
\end{enumerate}

Repeat the adjust step, watching the printed pot value, until it matches
the valve's known fully-open or fully-closed pot reading -- the PWM count
at that point is the calibration value to hardcode into production
firmware.

\subsection{Suggested Target Range}
\label{sec:suggested-range}

The range printed at each ``New Target'' prompt is computed at compile
time from a generic 5-turn-servo, 60:16-gear, 90\textdegree-travel ball
valve (the same geometry used across the production valves) -- it is a
starting-point hint, not a hard limit, since PWM counts still need
individual per-valve trimming on the bench:

\[
\text{start} \approx 2329 \qquad \text{end} \approx 3918
\]

These numbers come from the servo's 1100~Hz PWM frequency, its 500--2500~$\mu$s
mechanical range, the 60:16 gear ratio, and 90\textdegree\ of required valve
travel -- see \texttt{main.cpp}'s \texttt{constexpr} derivation if you need
to recompute them for a different gear ratio or travel angle.

%═══════════════════════════════════════════════════════════════════════════════
\section{Notes}
%═══════════════════════════════════════════════════════════════════════════════

\begin{itemize}[nosep]
  \item The Serial console has no timeout -- the tool waits indefinitely
        at every prompt.
  \item Backspace is supported while typing a value.
  \item There is no ``undo'' -- if a bad target damages the mechanism,
        power-cycle and re-enter values carefully.
\end{itemize}

%═══════════════════════════════════════════════════════════════════════════════
\section{Version History}
%═══════════════════════════════════════════════════════════════════════════════
\begin{center}
\begin{tabularx}{\textwidth}{@{}llY@{}}
\toprule
Version & Date & Notes \\
\midrule
0.1 & September 2026 & Initial user guide. \\
\bottomrule
\end{tabularx}
\end{center}

\end{document}
